\section{Relación entre $\theta_j$ y la frecuencia de rotación}

\begin{table}[H]
\centering
\caption{Índice del subespacio y características de frecuencia}
\begin{tabular}{cccc}
\toprule
\textbf{Posición del subespacio} & \textbf{Valor de $j$} & \textbf{$\theta_j$} & \textbf{Característica} \\
\midrule
Izquierda (índice bajo) & Pequeño & Grande (alta frecuencia) & Rotación rápida, codifica distancias cortas \\
Derecha (índice alto) & Grande & Pequeño (baja frecuencia) & Rotación lenta, codifica distancias largas \\
\bottomrule
\end{tabular}
\end{table}

\begin{knowledgebox}{Codificación multiescala de RoPE}
RoPE tiene distintas frecuencias de rotación en distintas dimensiones. Si se considera el vector de embedding como un vector fila:
\begin{itemize}
  \item Cuanto más a la izquierda esté el sub-space: mayor es la frecuencia, más sensible a la posición relativa
  \item Cuanto más a la derecha esté el sub-space: menor es la frecuencia, menos sensible a la posición relativa
\end{itemize}
Esta característica multiescala es la clave del éxito de RoPE.
\end{knowledgebox}

